\chapter*{Introduction to induction}
\addcontentsline{toc}{chapter}{Introduction to induction}
\markboth{Introduction to induction}{Introduction to induction}
\label{intro:induction}

A deductive science constructs a world from axioms and never has to look outside.
An inductive science works the other way round: it starts with the world as we
find it and asks which constructed world \emph{fits} it. Here we observe, and then
construct a model that matches what we observed.

\section*{Two meanings of \enquote{induction}}

The word is used for two unrelated things, and the difference matters here.
\emph{Mathematical induction} (\cref{thm:induction}) is a deductive proof method:
from $P(0)$ and $P(n) \Rightarrow P(n+1)$ it proves $P(n)$ for \emph{all} $n$,
with certainty. \emph{Scientific induction}, the subject of this part, goes from
finitely many observations to a general model. That step is never certain. No
number of white swans proves that all swans are white, and a single black one
refutes it.

\section*{The inductive cycle}

Every inductive science goes through the same cycle:
\begin{enumerate}
  \item \textbf{Observe.} Record what happens, together with how precisely it was
    measured.
  \item \textbf{Model.} Propose a mathematical structure that is meant to describe
    it. The model is itself a deductive object: a set with relations, maps and
    numbers, built with the tools of \cref{part:deduction}.
  \item \textbf{Deduce.} Work out, by proof or computation, what the model
    predicts, including things nobody has observed yet.
  \item \textbf{Compare.} Check the predictions against new observations. Keep the
    model if it fits, revise or replace it if it doesn't, and go back to step 1.
\end{enumerate}
Deduction is part of every round: step 3 is pure deduction. What makes the
science inductive is steps 1 and 4, where the world gets the last word.

\section*{What a model can claim}

A deductive theorem is true relative to its axioms, once and for all. An
inductive model has a weaker and different status:
\begin{itemize}
  \item It is never proved, only \emph{supported}: it has survived the comparisons
    made so far.
  \item It can be \emph{refuted} by a single reliable observation it can't
    accommodate. A model that no conceivable observation could contradict says
    nothing about the world.
  \item It has a \emph{domain of validity}. A model can be excellent at one scale
    and wrong at another, and still be the right tool inside its domain.
  \item \enquote{Fits} has to be made precise. Observations carry uncertainty, so
    comparing them with predictions needs \emph{probability} and
    \emph{statistics}. Probability theory is itself deductive. Statistical
    inference, which uses it to judge models against data, is the bridge between
    the two parts of this book.
\end{itemize}

\section*{One science, not many}

Physics, chemistry, biology and psychology are usually taught as separate
subjects. Here they are not separated, because they follow the same cycle and
differ only in the scale and kind of what they observe. A molecule is described
by a model built from models of atoms, a cell by models of molecules, behaviour by
models of nervous systems. So this part is ordered the way the deductive part is:
\emph{nothing is used before it has been introduced}. Here that applies to models
and to the observations supporting them, as well as to definitions.

\section*{Conventions for this part}

Each model will be presented with:
\begin{itemize}
  \item the \textbf{observations} it is meant to explain;
  \item its \textbf{assumptions}, stated as precisely as axioms;
  \item the \textbf{mathematics} it uses, with references to \cref{part:deduction};
  \item its \textbf{predictions}, and how they were tested;
  \item its known \textbf{limits}: where it fails, and which model takes over.
\end{itemize}

\section*{Planned chapters}

\begin{itemize}
  \item measurement, units and uncertainty;
  \item probability and statistical inference: how to judge a model against data;
  \item motion and forces;
  \item matter: from atoms to molecules and reactions;
  \item life: from molecules to cells, organisms and evolution;
  \item mind: from nervous systems to behaviour.
\end{itemize}
These chapters depend on mathematics that comes later in \cref{part:deduction}
(real numbers, calculus, linear algebra, probability), so they will be written once
that mathematics is in place.
